% GENERATED FILE. Edit canonical lessons, then run npm run generate:books.
% canonical-source-hash: fnv1a64:9c5a200d66d7c011
% canonical-lessons: TE-C476-sitosnasthiti, TE-C476-usnograta, TE-C476-tema, TE-C476-ukkapota, TE-C476-vedekku

\chapter{Talking About the Climate}
\label{ch:te-talking-about-the-climate}
\hlchaptermodality{476}

\begin{chapteropening}
\textbf{By the end of this chapter:} \emph{I can say the climate, temperature, humidity, muggy heat and to heat up.}

Complete the last lesson of chapter 476: I can say the climate, temperature, humidity, muggy heat and to heat up.
\end{chapteropening}

\section[\texorpdfstring{\te{శీతోష్ణస్థితి} (śītōṣṇasthiti)}{శీతోష్ణస్థితి (śītōṣṇasthiti)}]{\te{శీతోష్ణస్థితి} (śītōṣṇasthiti) — the climate}

\label{lesson:TE-C476-sitosnasthiti}



\begin{quote}
Before the new one: say the Telugu for irrigation, then the Telugu for a marsh.
\end{quote}

\subsection*{You'll want to know: \te{శీతోష్ణస్థితి}}
\textbf{\te{శీతోష్ణస్థితి}} — \emph{śītōṣṇasthiti} — \textquotedblleft{}the climate\textquotedblright{}.

\begin{quote}
\te{ఆ} \te{ప్రాంతం} \te{శీతోష్ణస్థితి} \te{వ్యవసాయానికి} \te{బాగుంటుంది}.
\end{quote}

\begin{quote}
\emph{ā prāntaṁ śītōṣṇasthiti vyavasāyāniki bāguṇṭundi.}
\end{quote}

\begin{quote}
The climate of that region is good for farming.
\end{quote}

\te{శీత}, cold, + \te{ఉష్ణ}, hot, + \te{స్థితి}, state: the state of cold and heat. In speech people often say \te{వాతావరణం}.

\begin{cousinweb}
\textbf{\te{శీతోష్ణస్థితి}} (\emph{śītōṣṇasthiti}) is built from three Sanskrit words: \emph{śīta}, \textquotedblleft{}cold\textquotedblright{} (CDIAL 12485), \emph{uṣṇa}, \textquotedblleft{}hot\textquotedblright{} (CDIAL 2389), and \emph{sthiti}, \textquotedblleft{}a state\textquotedblright{} (Monier-Williams).
\end{cousinweb}

\begin{grammarlens}[title={What you've built}]
One more word for getting around.
\end{grammarlens}

\subsection*{Your turn}
\begin{itemize}
\raggedright
  \item \emph{Say it:} \emph{śītōṣṇasthiti}
  \item \emph{Say it:} \emph{śītōṣṇasthiti}, once more
  \item \emph{Say it:} \emph{ā prāntaṁ śītōṣṇasthiti vyavasāyāniki bāguṇṭundi.} — The climate of that region is good for farming.
  \item \emph{Recall:} say the Telugu for irrigation, then the Telugu for a marsh, then say \emph{śītōṣṇasthiti} again
\end{itemize}

\begin{tcolorbox}[breakable,colback=teal!4,colframe=teal!35!black,title={Before you move on}]
What does \te{శీతోష్ణస్థితి} mean? (\textquotedblleft{}The climate\textquotedblright{}.) Say it once more.
\end{tcolorbox}
\section[\texorpdfstring{\te{ఉష్ణోగ్రత} (uṣṇōgrata)}{ఉష్ణోగ్రత (uṣṇōgrata)}]{\te{ఉష్ణోగ్రత} (uṣṇōgrata) — temperature}

\label{lesson:TE-C476-usnograta}



\begin{quote}
Before the new one: say the Telugu for a marsh, then the Telugu for the climate.
\end{quote}

\subsection*{You'll want to know: \te{ఉష్ణోగ్రత}}
\textbf{\te{ఉష్ణోగ్రత}} — \emph{uṣṇōgrata} — \textquotedblleft{}temperature\textquotedblright{}.

\begin{quote}
\te{ఇవాళ} \te{ఉష్ణోగ్రత} \te{నలభై} \te{డిగ్రీలు} \te{దాటింది}.
\end{quote}

\begin{quote}
\emph{ivāḷa uṣṇōgrata nalabhai ḍigrīlu dāṭindi.}
\end{quote}

\begin{quote}
Today the temperature went above forty degrees.
\end{quote}

\te{ఉష్ణ}, heat, + \te{ఉగ్రత}, fierceness: how fierce the heat is.

\begin{cousinweb}
\textbf{\te{ఉష్ణోగ్రత}} (\emph{uṣṇōgrata}) is built from Sanskrit \emph{uṣṇa}, \textquotedblleft{}hot, heat\textquotedblright{} (CDIAL 2389), and \emph{ugratā}, \textquotedblleft{}fierceness\textquotedblright{} (Monier-Williams).
\end{cousinweb}

\begin{grammarlens}[title={What you've built}]
One more word for getting around.
\end{grammarlens}

\subsection*{Your turn}
\begin{itemize}
\raggedright
  \item \emph{Say it:} \emph{uṣṇōgrata}
  \item \emph{Say it:} \emph{uṣṇōgrata}, once more
  \item \emph{Say it:} \emph{ivāḷa uṣṇōgrata nalabhai ḍigrīlu dāṭindi.} — Today the temperature went above forty degrees.
  \item \emph{Recall:} say the Telugu for a marsh, then the Telugu for the climate, then say \emph{uṣṇōgrata} again
\end{itemize}

\begin{tcolorbox}[breakable,colback=teal!4,colframe=teal!35!black,title={Before you move on}]
What does \te{ఉష్ణోగ్రత} mean? (\textquotedblleft{}Temperature\textquotedblright{}.) Say it once more.
\end{tcolorbox}
\section[\texorpdfstring{\te{తేమ} (tēma)}{తేమ (tēma)}]{\te{తేమ} (tēma) — humidity; moisture}

\label{lesson:TE-C476-tema}



\begin{quote}
Before the new one: say the Telugu for the climate, then the Telugu for temperature.
\end{quote}

\subsection*{You'll want to know: \te{తేమ}}
\textbf{\te{తేమ}} — \emph{tēma} — \textquotedblleft{}humidity; moisture\textquotedblright{}.

\begin{quote}
\te{సముద్రం} \te{దగ్గర} \te{గాలిలో} \te{తేమ} \te{ఎక్కువ}.
\end{quote}

\begin{quote}
\emph{samudraṁ daggara gālilō tēma ekkuva.}
\end{quote}

\begin{quote}
Near the sea there is a lot of moisture in the air.
\end{quote}

\te{తేమగా} \te{ఉంది}, it is damp.

\begin{cousinweb}
\textbf{\te{తేమ}} (\emph{tēma}) is a Telugu word for damp and moisture whose origin is not settled.
\end{cousinweb}

\begin{grammarlens}[title={What you've built}]
One more word for getting around.
\end{grammarlens}

\subsection*{Your turn}
\begin{itemize}
\raggedright
  \item \emph{Say it:} \emph{tēma}
  \item \emph{Say it:} \emph{tēma}, once more
  \item \emph{Say it:} \emph{samudraṁ daggara gālilō tēma ekkuva.} — Near the sea there is a lot of moisture in the air.
  \item \emph{Recall:} say the Telugu for the climate, then the Telugu for temperature, then say \emph{tēma} again
\end{itemize}

\begin{tcolorbox}[breakable,colback=teal!4,colframe=teal!35!black,title={Before you move on}]
What does \te{తేమ} mean? (\textquotedblleft{}Humidity; moisture\textquotedblright{}.) Say it once more.
\end{tcolorbox}
\section[\texorpdfstring{\te{ఉక్కపోత} (ukkapōta)}{ఉక్కపోత (ukkapōta)}]{\te{ఉక్కపోత} (ukkapōta) — muggy heat}

\label{lesson:TE-C476-ukkapota}



\begin{quote}
Before the new one: say the Telugu for temperature, then the Telugu for humidity.
\end{quote}

\subsection*{You'll want to know: \te{ఉక్కపోత}}
\textbf{\te{ఉక్కపోత}} — \emph{ukkapōta} — \textquotedblleft{}muggy heat\textquotedblright{}.

\begin{quote}
\te{వర్షం} \te{ముందు} \te{ఉక్కపోత} \te{భరించలేకపోయాం}.
\end{quote}

\begin{quote}
\emph{varṣaṁ mundu ukkapōta bhariñcalēkapōyāṁ.}
\end{quote}

\begin{quote}
Before the rain we could not bear the muggy heat.
\end{quote}

\te{ఉక్క}, closeness of the air, + \te{పోత}. The heat that makes you sweat without any breeze.

\begin{cousinweb}
\textbf{\te{ఉక్కపోత}} (\emph{ukkapōta}) is a Telugu compound: \te{ఉక్క}, \textquotedblleft{}sultriness, closeness of the weather\textquotedblright{} (DEDR 568), native Dravidian, and -\te{పోత}, whose origin is not settled.
\end{cousinweb}

\begin{grammarlens}[title={What you've built}]
One more word for getting around.
\end{grammarlens}

\subsection*{Your turn}
\begin{itemize}
\raggedright
  \item \emph{Say it:} \emph{ukkapōta}
  \item \emph{Say it:} \emph{ukkapōta}, once more
  \item \emph{Say it:} \emph{varṣaṁ mundu ukkapōta bhariñcalēkapōyāṁ.} — Before the rain we could not bear the muggy heat.
  \item \emph{Recall:} say the Telugu for temperature, then the Telugu for humidity, then say \emph{ukkapōta} again
\end{itemize}

\begin{tcolorbox}[breakable,colback=teal!4,colframe=teal!35!black,title={Before you move on}]
What does \te{ఉక్కపోత} mean? (\textquotedblleft{}Muggy heat\textquotedblright{}.) Say it once more.
\end{tcolorbox}
\section[\texorpdfstring{\te{వేడెక్కు} (vēḍekku)}{వేడెక్కు (vēḍekku)}]{\te{వేడెక్కు} (vēḍekku) — to heat up; to grow hot}

\label{lesson:TE-C476-vedekku}



\begin{quote}
Before the new one: say the Telugu for humidity, then the Telugu for muggy heat.
\end{quote}

\subsection*{You'll want to know: \te{వేడెక్కు}}
\textbf{\te{వేడెక్కు}} — \emph{vēḍekku} — \textquotedblleft{}to heat up; to grow hot\textquotedblright{}.

\begin{quote}
\te{మధ్యాహ్నానికి} \te{గది} \te{బాగా} \te{వేడెక్కింది}.
\end{quote}

\begin{quote}
\emph{madhyāhnāniki gadi bāgā vēḍekkindi.}
\end{quote}

\begin{quote}
By midday the room had got very hot.
\end{quote}

Past (it …): \te{వేడెక్కింది}, \emph{vēḍekkindi}.

\te{వేడి}, heat, + \te{ఎక్కు}, to climb: heat climbing. Also said of a heated argument. Its subject is a thing, so the past shown is the third person: \te{వేడెక్కింది}, it got hot.

\begin{cousinweb}
\textbf{\te{వేడెక్కు}} (\emph{vēḍekku}) is native Dravidian: \te{వేడి}, \textquotedblleft{}heat\textquotedblright{} (DEDR 5517), and \te{ఎక్కు}, \textquotedblleft{}to climb, to rise\textquotedblright{} (DEDR 766).
\end{cousinweb}

\begin{grammarlens}[title={What you've built}]
One more word for getting around.
\end{grammarlens}

\subsection*{Your turn}
\begin{itemize}
\raggedright
  \item \emph{Say it:} \emph{vēḍekku}
  \item \emph{Say it:} \emph{vēḍekku}, once more
  \item \emph{Say it:} \emph{madhyāhnāniki gadi bāgā vēḍekkindi.} — By midday the room had got very hot.
  \item \emph{Recall:} say the Telugu for humidity, then the Telugu for muggy heat, then say \emph{vēḍekku} again
\end{itemize}

\begin{tcolorbox}[breakable,colback=teal!4,colframe=teal!35!black,title={Before you move on}]
What does \te{వేడెక్కు} mean? (\textquotedblleft{}To heat up; to grow hot\textquotedblright{}.) Say it once more.
\end{tcolorbox}
